% s-232-p232-m002; original PDF 232, 233
\Aphorism{s-232-p232-m002}{327}
Ich möchte gern zur „Gerechtigkeit“ gelangen, nicht im frömmelnd-passiven Sinne, sondern als gerechter Richter und Betrachter der Dinge, vor dem das Einzelne mit seiner heftigen Anmassung und leidenschaftlichen Beschränktheit zur Ruhe verwiesen und ein friedliches Nebeneinander der widersprechenden Elemente möglich wird. Keine völlige Objectivität, aber doch ein Freisein von den gröberen Dünsten und Trübheiten des Subjects; weder Einsiedler-Schweigen noch Majoritätenlärm, sondern der consensus einer Minderheit billig denkender und weitblickender Geister. Der bloss nachschwatzende Heerdenmensch wäre von dieser Gesellschaft ebenso ausgeschlossen wie der fanatische Einzelne und Revolutionär, der die „Wahrheit“ sagt und wenn die Welt in Stücken ginge. Auf dem Wege zu diesem Hochplateau, das weder Niederung noch isolirte Spitze sein will, ist viel „Individuum“ zu überwinden, viel bornirter Enthusiasmus und sancta simplicitas, viel Apostel- und Heilandswahn mit feierlicher Verketzerung der „blinden Heiden“ — ja sogar viel gedankenlose Unart im Urtheilen und Sprechen. Die Sprache, das uralte Werkzeug der Ungerechtigkeit, hat um uns moderne Menschen einen Ballast von willkürlichen und anmasslichen \Foreign{façons de parler} gehäuft, von wüthender Parteilichkeit und schnöder Ironie, von kategorischem Grossmaul und unverschämtem Superlativismus: es wimmelt in allen Sprachen von entschieden und unbedingt, von absolument und senza dubbio — nichts geht dem modernen Sprachton mehr wider die Gewöhnung als ein wenig Ruhe, Bescheidenheit, Skepsis. Wer sich heute vorsichtig und gerecht ausdrücken will, hat einen förmlichen Kampf mit der barbarischen „Schneidigkeit“ und Autoritätswuth unserer Parteihäuptlings-Fanfarensprache zu bestehen; er wird in den gewundenen umschweifenden ausweichenden Leisetreterstil verfallen, der so viele wissenschaftliche Arbeiten Derer, die nach „Gerechtigkeit“ trachten, unlesbar macht. Wenn \Sperr{Burckhardt’s} Cultur der Renaissance für nichts anderes zu bewundern wäre, so allein schon für ihre Goethische Art, zugleich klar, kurz und massvoll im Ausdruck ihrer Überzeugung zu sein; sie trägt antike Sandalen und kommt weder mit Sporen geklirrt noch auf Strümpfen geschlichen. — Wenn es dem Autor anstünde, von seinem eigenen Wollen und Verfehlen zu reden, so wäre hier der Ort, für die vielfache Ungerechtigkeit dieses Buches (die nicht immer dem Autor, sondern bisweilen dem entraînement durch die moderne Sprache zur Last fällt) um Verzeihung zu bitten.
\par

